% ==============================================================================
% KATA PENGANTAR
% Sesuai Format Skripsi Kating (FST UNIBA Madura)
% ==============================================================================

\clearpage
\phantomsection
\pagestyle{frontstyle}
\addcontentsline{toc}{chapter}{KATA PENGANTAR}

\begin{center}
    {\fontsize{14pt}{18pt}\selectfont \bfseries KATA PENGANTAR}
\end{center}

\vspace{0.3cm}

\textit{Alhamdulillahi Rabbil 'Alamin}, segala puji dan syukur penulis panjatkan kehadirat Allah SWT atas limpahan berkah, rahmat, dan karunia-Nya, sehingga penulis dapat menyelesaikan penulisan proposal skripsi yang berjudul \textbf{“\skripsiJudul”} ini sebagai salah satu syarat untuk memperoleh gelar \skripsiGelar\ (\skripsiSingkatanGelar) pada Program Studi \skripsiProdi, \skripsiFakultas, \skripsiUniversitas.

Penyusunan skripsi ini tidak lepas dari bimbingan, arahan, dan dukungan dari berbagai pihak. Oleh karena itu, pada kesempatan ini penulis menyampaikan terima kasih yang tulus kepada:

\begin{enumerate}\singlespacing\setlength{\itemsep}{1.5pt}\setlength{\parsep}{0pt}\setlength{\topsep}{2pt}
    \item Bapak \skripsiDekan, selaku Dekan Fakultas Sains dan Teknologi Universitas K.H. Bahaudin Mudhary Madura.
    \item Bapak Mohammad Iqbal Bachtiar, S.T., M.T., selaku Wakil Dekan Fakultas Sains dan Teknologi Universitas K.H. Bahaudin Mudhary Madura.
    \item Bapak \skripsiKaprodi, selaku Ketua Program Studi \skripsiProdi\ Universitas K.H. Bahaudin Mudhary Madura.
    \item Bapak \skripsiPembimbingSatu, selaku Dosen Pembimbing I atas segala bimbingan, telaah ilmiah, dan arahannya.
    \item Bapak \skripsiPembimbingDua, selaku Dosen Pembimbing II atas bimbingan teknis, saran konstruktif, dan motivasinya.
    \item Seluruh Dosen dan Tenaga Kependidikan Program Studi \skripsiProdi\ Fakultas Sains dan Teknologi.
    \item Teruntuk Kedua Orang Tua tercinta, Ayahanda dan Ibunda, serta segenap keluarga besar atas doa tulus dan pengorbanan tanpa pamrih.
    \item Rekan-rekan mahasiswa seperjuangan Program Studi \skripsiProdi\ atas kerja sama dan kebersamaannya.
\end{enumerate}

Penulis menyadari bahwa dalam penulisan naskah ini masih terdapat keterbatasan. Kritik dan saran membangun sangat penulis harapkan demi penyempurnaan karya ilmiah ini. Semoga skripsi ini bermanfaat bagi kita semua. Aamiin.

\vspace{0.3cm}

\begin{flushright}
\begin{tabular}{l}
    \skripsiKota, \dots\ \dots\ \skripsiTahun\\[1.4cm]
    \skripsiPenulis
\end{tabular}
\end{flushright}
